\section{Some Applications of Integral Calculus}

\begin{theorem}[Irrationality of $\pi$]
\mymargin{$\pi$ is irrational}
\index{$\pi$!irrationality}
\index{irrationality!of $\pi$}
The number $\pi$ is irrational.
\end{theorem}
%
\begin{proof}
Our first objective is to find recursive formulas for the computation of the
following integral:
\[
  I_n = \int_0^\pi x^n (\pi-x)^n \sin x\, dx.
\]
First, observe that we have
\[
  I_0 = \int_0^\pi \sin x\, dx = \Enclose{-\cos x}_0^\pi = 2
\]
and, using integration by parts:
\begin{align*}
  I_1 &= \int_0^\pi x(\pi-x)\sin x\, dx\\
      &= \Enclose{x(\pi-x)(-\cos x)}_0^\pi
       -\int_0^\pi (\pi-2x)(-\cos x)\, dx \\
      &= \int_0^\pi (\pi-2x)\cos x\, dx
      = \pi \int_0^\pi \cos x\, dx - 2 \int_0^\pi x\cos x\, dx \\
      &= - 2 \enclose{\Enclose{x\sin x}_0^\pi - \int_0^\pi \sin x\, dx}
      = 2 \Enclose{\cos x}_0^\pi = 4.
\end{align*}
For $n\ge 2$, we proceed similarly. The function $x^n (\pi-x)^n$ vanishes at the
integration limits, so by integrating by parts, we obtain:
\[
    I_n = \int_0^\pi \Enclose{n x^{n-1}(\pi-x)^n - n x^n (\pi-x)^{n-1}} \cos x\,
  dx
\]
The function in the square brackets also vanishes at the integration limits, and
thus integrating again by parts, we obtain:
\begin{align*}
  I_n &= -\int_0^\pi \big[n(n-1) x^{n-2}(\pi-x)^n
    - 2n^2 x^{n-1}(\pi-x)^{n-1} \\
  &\qquad\quad   + n(n-1) x^n(\pi-x)^{n-2}\big] \sin(x)\, dx \\
    &= 2n^2 I_{n-1} - (n^2-n) \int_0^\pi \Enclose{(\pi-x)^2 +
  x^2}x^{n-2}(\pi-x)^{n-2} \sin x\, dx \\
    &= 2n^2 I_{n-1} - (n^2-n) \int_0^\pi \Enclose{\pi^2 -
  2x(\pi-x)}x^{n-2}(\pi-x)^{n-2}\sin x\, dx \\
  &= 2n^2 I_{n-1} - (n^2-n) \pi^2 I_{n-2} + 2(n^2-n) I_{n-1}
\end{align*}
from which
\begin{align}\label{eq:9530978}
   I_n = (4n^2-2n)I_{n-1} -(n^2-n)\pi^2 I_{n-2}.
\end{align}

Now, suppose for contradiction that $\pi = \frac p q$ with $p,q\in \NN$ and
consider the sequence
\[
   a_n = \frac{q^{2n}}{n!} I_n.
\]
From the relation~\eqref{eq:9530978}, we deduce a recursive relation for $a_n$:
\begin{align*}
    a_n &= \frac{q^{2n}}{n!}\enclose{-(n^2-n)\frac{p^2}{q^2}
  \frac{(n-2)!}{q^{2n-4}}a_{n-2}
  + (4n^2-2n) \frac{(n-1)!}{q^{2n-2}}a_{n-1}} \\
  &= - q^2 p^2 a_{n-2} + (4n-2)q^2 a_{n-1}
\end{align*}
which tells us, in particular, that if $a_{n-2}$ and $a_{n-1}$ are integers,
then $a_n$ is also an integer (since it is a sum of products of integers).
Since the first two terms:
\[
  a_0 = I_0 = 2, \qquad
  a_1 = q^2 I_1 = 4 q^2
\]
are integers, we can deduce by induction that all terms of the sequence $a_n$
are integers.

On the other hand, observing that $y=x(\pi-x)$ is the graph of a parabola with
axis $x=\frac \pi 2$, we know that for every $x$, $x(\pi-x) \le
\frac{\pi^2}{4}$, and thus
\[
  0 \le x^n(\pi-x)^n \sin x
    \le x^n (\pi-x)^n \le \frac{\pi^{2n}}{4^n}.
\]
Therefore,
\[
  0 \le I_n \le \frac{\pi^{2n+1}}{4^n}
\]
which implies
\[
  0 \le a_n \le \frac{q^{2n}}{n!}\frac{\pi^{2n+1}}{4^n}.
\]
Since $n! \gg C^n$, we can thus assert that $a_n\to 0$.
On the other hand, we have seen that $a_n \in \ZZ$ and certainly $a_n > 0$ since
$I_n \neq 0$ because $f_n$ is a continuous, non-negative, and non-zero function.
But it is not possible for a sequence of positive integers to converge to zero:
we have thus reached a contradiction.
\end{proof}

\begin{theorem}[Wallis Product]
\label{th:wallis}%
\mymark{*}%
\mymargin{Wallis product}%
\index{$\pi$!Wallis product}%
\index{Wallis!approximation of $\pi$}%
\index{formula!of Wallis}%
We have
\[
  \frac{\pi}{2}
  = \prod_{k=1}^{+\infty} \frac{(2k)^2}{(2k-1)(2k+1)}
  = \frac{2}{1}
  \cdot \frac{2}{3}
  \cdot \frac{4}{3}
  \cdot \frac{4}{5}
  \cdot \frac{6}{5}
  \cdot \frac{6}{7}
  \cdot \frac{8}{7}
  \cdot \frac{8}{9}
  \dots
\]
And consequently, we obtain the following asymptotic estimate for the
\emph{central binomial coefficient}%
\mymargin{central binomial}%
\index{binomial!central}
\[
 {2n \choose n} \sim \frac{4^n}{\sqrt{\pi n}}
 \qquad \text{as $n\to +\infty$.}
\]
\end{theorem}
%
\begin{proof}
Consider the sequence of integrals:
\[
  I_n = \int_0^\pi \sin^n(x)\, dx.
\]
Since $0\le \sin^n(x) \le 1$ for every $x\in [0,\pi]$ and since
$\sin^{n+1}(x)\le \sin^n(x)$ for every $x\in[0,\pi]$, it is clear that $I_n$ is
a decreasing sequence of positive numbers.

From a direct calculation, we find that
\[
  I_0 = \int_0^\pi 1\, dx = \pi, \qquad
  I_1 = \int_0^\pi \sin(x)\, dx = \Enclose {-\cos x}_0^\pi = 2.
\]
If $n\ge 2$,
integrating by parts, we find instead a recursive relation
\begin{align*}
 I_n &= \int_0^\pi \sin^{n-1}(x) \sin x\, dx \\
          &= \Enclose{-\sin^{n-1} x \cos x }_0^\pi + (n-1)\int_0^\pi \sin^{n-2}
     x \cos^2 x \, dx \\
     &= 0 + (n-1)\int_0^\pi\sin^{n-2}x\cdot (1-\sin^2 x)\, dx \\
     &= (n-1) I_{n-2} - (n-1) I_{n}
\end{align*}
from which
\begin{equation}
\label{eq:187464}%
  I_n = \frac{n-1}{n} I_{n-2}.
\end{equation}
This recursive formula allows us to calculate separately the even and odd terms
of the sequence using the notation of the \emph{semi-factorial} (see
Remark~\ref{rem:doppio_fattoriale}):
\begin{align*}
  I_{2n}
    &= \frac{2n-1}{2n}\cdot \frac{2n-3}{2n-2} \cdots \frac {3}{4}\cdot
  \frac{1}{2} \cdot I_0
  = \frac{(2n-1)!!}{(2n)!!} \cdot \pi \\
  I_{2n+1} &= \frac{2n}{2n+1}\cdot \frac {2n-2}{2n-1} \cdots
  \frac{4}{5}\cdot \frac{2}{3} \cdot I_1
  = \frac{(2n)!!}{(2n+1)!!}\cdot 2.
\end{align*}

Recalling that $I_n$ is decreasing, we observe that
\[
  \frac{I_{2n+2}}{I_{2n}}
  \le \frac{I_{2n+1}}{I_{2n}}
  \le \frac{I_{2n}}{I_{2n}} = 1
\]
and thanks to~\eqref{eq:187464}
\[
 \frac{I_{2n+2}}{I_{2n}}
 = \frac{2n+1}{2n+2} \to 1
\]
from which we obtain that $\frac{I_{2n+1}}{I_{2n}}\to 1$.

In conclusion, since $\pi$ appears in the formula for the even terms but not in
that for the odd terms, and since the two expressions are asymptotic, we can
obtain the formula for calculating $\frac \pi 2$:
\begin{align*}
\frac{\pi}{2}
&= \frac{\pi}{2} \lim_{n\to +\infty} \frac{I_{2n+1}}{I_{2n}}
= \lim_{n\to +\infty} \frac{\frac{(2n)!!}{(2n+1)!!}}{\frac{(2n-1)!!}{(2n)!!}}
= \lim_{n\to +\infty} \frac{\enclose{(2n)!!}^2}{(2n+1)!!(2n-1)!!} \\
&= \lim_{n\to +\infty} \frac{(2n)(2n)(2n-2)(2n-2) \cdots 4\cdot 4 \cdot 2 \cdot
2}
  {(2n+1)(2n-1)(2n-1)(2n-3)(2n-3)\cdots 3 \cdot 3 \cdot 1}.
\end{align*}

To obtain an asymptotic estimate of the central binomial coefficient, recall
that (Remark~\ref{rem:doppio_fattoriale})
\begin{align*}
  (2n)!! &= 2^n n! \\
  (2n+1)!! &= \frac{(2n+1)!}{2^n n!}
\end{align*}
thus the asymptotic estimate of Wallis can be written as
\begin{align*}
 \sqrt{\frac{\pi}{2}}
 &\sim \frac{(2n)!!}{\sqrt{2n+1}(2n-1)!!}
  = \frac{2^n n!\sqrt{2n+1}}{\frac{(2n+1)!}{2^n n!}} \\
 &= \frac{\enclose{2^n (n!)}^2 \sqrt{2n+1}}{(2n+1)!} 
 \sim \frac{4^n (n!)^2}{\sqrt{2n}(2n)!}
\end{align*}
from which
\[
{2n \choose n}
= \frac{(2n)!}{(n!)^2}
\sim \frac{4^n}{\sqrt{\pi n}}.
\]
\end{proof}

\begin{theorem}[Stirling's Formula]
\label{th:stirling}%
\mymark{*}%
\mymargin{Stirling's formula}%
\index{formula!of Stirling}%
\index{Stirling formula of}%
\index{factorial!Stirling's formula}%
We have
\[
  n! \sim \sqrt{2\pi n}\cdot \frac{n^n}{e^n}
  \qquad \text{as $n\to +\infty$.}
\]
\end{theorem}
%
\begin{proof}
Observe that the thesis:
\[
  \lim_{n\to +\infty}\frac{\sqrt{n} \cdot n^n}{n!\cdot e^n } = \frac{1}{\sqrt{2\pi}}
\]
is equivalent, by taking logarithms, to proving that
\[
   \lim_{n\to +\infty} \enclose{\frac 1 2 \ln n + n \ln n - \ln(n!) - n}
   = -\frac 1 2 \ln (2\pi).
\]
The first part of the proof will aim to demonstrate that the preceding limit
exists and is finite. Finally, to find the exact value, we will use Wallis's
formula.
The quantity $n\ln n - n$ is, up to an additive constant, the integral of the
logarithm over the interval $[1,n]$.
We have already seen in Example~\ref{ex:498124} that $\ln(n!)$ is obtained by
approximating this integral with rectangles.
If we proceed instead with an approximation using trapezoids, we obtain a more
precise estimate that will give us the additional term $\frac 1 2 \ln n$,
necessary for the limit to converge.

Observe in general that if $f:[a,b]\to \RR$ is a concave function, then the
graph of $f$ is bounded between the line passing through the endpoints
$(a,f(a))$, $(b,f(b))$ (secant line) and any tangent line, for example, the
tangent line at $((a+b)/2,f((a+b)/2))$.
Consequently, the area under the graph, i.e., $\int_a^b f(x)\, dx$, is bounded
between the areas of the two corresponding right trapezoids. The height of both
trapezoids is $(b-a)$, and the area is calculated by multiplying the height by
the average of the bases. In the case of the trapezoid with the oblique side on
the secant, the average of the bases is $(f(a)+f(b))/2$, in the case of the
trapezoid with the oblique side on the tangent line at the midpoint, the average
of the bases is equal to the section at the midpoint: $f((a+b)/2)$. Thus, for
every concave $f$:
\[
   (b-a)\cdot  \frac{f(a) + f(b)}{2}
   \le \int_a^b f(x) \, dx
   \le (b-a) \cdot f\enclose{\frac{a+b}{2}}.
\]

Applying these estimates to the function $f(x) = \ln x$ on the interval $[k,
k+1]$, we obtain:
\[
\frac{\ln(k) + \ln(k+1)}{2}
\le \int_k^{k+1} \ln x\, dx
\le \ln \enclose{k+\frac 1 2}.
\]

Let $a_k$ be the difference between the first two quantities.
Clearly, $a_k\ge 0$
and it results in
\begin{align*}
a_k
 &= \int_k^{k+1} \ln x\, dx - \frac{\ln (k) + \ln(k+1)}{2}\\
 & \le \ln\enclose{k+\frac 1 2} - \frac{\ln(k) + \ln(k+1)}{2} \\
 & = \frac 1 2 \ln\frac{\enclose{k+\frac 1 2}^2}{k(k+1)}
  = \frac 1 2 \ln\enclose{1+\frac{1}{4(k^2+k)}}\\
 & = \frac 1 {8k^2} + o\enclose{\frac 1 {k^2}}
 \qquad \text{as $k\to+\infty$.}
\end{align*}
The series $\sum a_k$ is thus convergent, meaning the sequence
\[
  A_n = \sum_{k=1}^n a_k
\]
is convergent. Let $\ell$ be the limit of $A_n$. We have:
\begin{align*}
A_n
&= \sum_{k=1}^{n-1} \int_{k}^{k+1} \ln x\, dx -
\sum_{k=1}^{n-1} \frac{\ln k + \ln(k+1)}{2} \\
&= \int_{1}^{n} \ln x\, dx - \sum_{k=2}^{n-1} \ln k - \frac{\ln 1}{2}
-\frac{\ln n}{2} \\
&= \Enclose{ x \ln x -x}_1^n
- \sum_{k=1}^n \ln k + \frac{\ln n}{2} \\
&= n \ln n - n + 1 - \ln(n!) + \frac {\ln n}{2} \to \ell
 \end{align*}
from which
\[
e^{A_n} = \frac{n^n\cdot e \cdot \sqrt{n}}{e^n \cdot n!} \to e^\ell
\]
hence, setting $c=e/e^\ell$,
\[
  n! \sim c\cdot\frac{n^n \sqrt{n}}{e^n}.
\]
To determine the unknown constant $c$, we can exploit Wallis's formula on the
central binomial coefficient:
\begin{align*}
\frac{4^n}{\sqrt{\pi n}}
 &\sim
{2n \choose n}
= \frac{(2n)!}{(n!)^2}
\sim \frac{\displaystyle c \frac{(2n)^{2n}\sqrt{2n}}{e^{2n}}}
{\displaystyle\enclose{c\frac{n^n\sqrt n}{e^n}}^2} \\
&= \frac{2^{2n}\sqrt{2}}{c\sqrt{n}}
= \frac{4^n \sqrt 2}{c\sqrt n}
\end{align*}
from which we obtain
\[
  \frac{1}{\sqrt \pi} \sim \frac{\sqrt 2}{c}
\]
and therefore $c=\sqrt{2\pi}$.
\end{proof}